% Options for packages loaded elsewhere
\PassOptionsToPackage{unicode}{hyperref}
\PassOptionsToPackage{hyphens}{url}
%
\documentclass[
  11pt,
]{article}
\usepackage{amsmath,amssymb}
\usepackage{iftex}
\ifPDFTeX
  \usepackage[T1]{fontenc}
  \usepackage[utf8]{inputenc}
  \usepackage{textcomp} % provide euro and other symbols
\else % if luatex or xetex
  \usepackage{unicode-math} % this also loads fontspec
  \defaultfontfeatures{Scale=MatchLowercase}
  \defaultfontfeatures[\rmfamily]{Ligatures=TeX,Scale=1}
\fi
\usepackage{lmodern}
\ifPDFTeX\else
  % xetex/luatex font selection
\fi
% Use upquote if available, for straight quotes in verbatim environments
\IfFileExists{upquote.sty}{\usepackage{upquote}}{}
\IfFileExists{microtype.sty}{% use microtype if available
  \usepackage[]{microtype}
  \UseMicrotypeSet[protrusion]{basicmath} % disable protrusion for tt fonts
}{}
\makeatletter
\@ifundefined{KOMAClassName}{% if non-KOMA class
  \IfFileExists{parskip.sty}{%
    \usepackage{parskip}
  }{% else
    \setlength{\parindent}{0pt}
    \setlength{\parskip}{6pt plus 2pt minus 1pt}}
}{% if KOMA class
  \KOMAoptions{parskip=half}}
\makeatother
\usepackage{xcolor}
\usepackage[margin=2.2cm]{geometry}
\usepackage{longtable,booktabs,array}
\usepackage{calc} % for calculating minipage widths
% Correct order of tables after \paragraph or \subparagraph
\usepackage{etoolbox}
\makeatletter
\patchcmd\longtable{\par}{\if@noskipsec\mbox{}\fi\par}{}{}
\makeatother
% Allow footnotes in longtable head/foot
\IfFileExists{footnotehyper.sty}{\usepackage{footnotehyper}}{\usepackage{footnote}}
\makesavenoteenv{longtable}
\setlength{\emergencystretch}{3em} % prevent overfull lines
\providecommand{\tightlist}{%
  \setlength{\itemsep}{0pt}\setlength{\parskip}{0pt}}
\setcounter{secnumdepth}{-\maxdimen} % remove section numbering
\ifLuaTeX
  \usepackage{selnolig}  % disable illegal ligatures
\fi
\IfFileExists{bookmark.sty}{\usepackage{bookmark}}{\usepackage{hyperref}}
\IfFileExists{xurl.sty}{\usepackage{xurl}}{} % add URL line breaks if available
\urlstyle{same}
\hypersetup{
  pdftitle={Worksheet 4D: Transcendental Functions via Series},
  hidelinks,
  pdfcreator={LaTeX via pandoc}}

\title{Worksheet 4D: Transcendental Functions via Series}
\author{}
\date{}

\begin{document}
\maketitle

\begin{quote}
\textbf{What This Worksheet Is About}

So far, we've worked with polynomials. But what about functions like
\(e^x\), \(\sin(x)\), \(\cos(x)\), and \(\ln(x)\)?

Here's the beautiful secret: \textbf{These functions can be written as
infinite polynomials (series)!}

Once we write them as polynomials, we can differentiate and integrate
term by term using the power rule. This lets us \emph{discover} their
derivatives rather than just memorizing them.
\end{quote}

\begin{center}\rule{0.5\linewidth}{0.5pt}\end{center}

\hypertarget{part-a-the-exponential-function-ex-18-problems}{%
\subsection{\texorpdfstring{Part A: The Exponential Function \(e^x\) (18
problems)}{Part A: The Exponential Function e\^{}x (18 problems)}}\label{part-a-the-exponential-function-ex-18-problems}}

\hypertarget{the-series-definition}{%
\subsubsection{The Series Definition}\label{the-series-definition}}

The number \(e \approx 2.71828...\) is defined as:

\[e = 1 + 1 + \frac{1}{2!} + \frac{1}{3!} + \frac{1}{4!} + \cdots\]

And the exponential function \(e^x\) is:

\[\boxed{e^x = 1 + x + \frac{x^2}{2!} + \frac{x^3}{3!} + \frac{x^4}{4!} + \cdots = \sum_{n=0}^{\infty} \frac{x^n}{n!}}\]

\textbf{Reminder:}
\(n! = n \times (n-1) \times \cdots \times 2 \times 1\), and \(0! = 1\)
by convention.

\begin{center}\rule{0.5\linewidth}{0.5pt}\end{center}

\textbf{Understanding the Series:}

\textbf{1.} Calculate the first five terms of the series for
\(e^1 = e\):

\(1 + 1 + \frac{1}{2} + \frac{1}{6} + \frac{1}{24} ={}\)
\_\_\_\_\_\_\_\_\_\_

\textbf{2.} Compare your answer to \(e \approx 2.71828\). How close?

\textbf{3.} Write out the first 6 terms of the series for \(e^x\):

\(e^x = 1 + x +{}\) \_\_\_\_\_\_\_\_\_\_ \(+\) \_\_\_\_\_\_\_\_\_\_
\(+\) \_\_\_\_\_\_\_\_\_\_ \(+\) \_\_\_\_\_\_\_\_\_\_ \(+ \cdots\)

\textbf{4.} Use the first 4 terms to approximate \(e^{0.5}\):

\(e^{0.5} \approx 1 + 0.5 + \frac{0.25}{2} + \frac{0.125}{6} ={}\)
\_\_\_\_\_\_\_\_\_\_

Actual value: \(e^{0.5} \approx 1.6487\). How close?

\begin{center}\rule{0.5\linewidth}{0.5pt}\end{center}

\hypertarget{differentiating-ex-term-by-term}{%
\subsubsection{\texorpdfstring{Differentiating \(e^x\) Term by
Term}{Differentiating e\^{}x Term by Term}}\label{differentiating-ex-term-by-term}}

\textbf{5.} Differentiate each term of the series
\(e^x = 1 + x + \frac{x^2}{2!} + \frac{x^3}{3!} + \frac{x^4}{4!} + \cdots\)

\begin{longtable}[]{@{}ll@{}}
\toprule\noalign{}
Term & Derivative \\
\midrule\noalign{}
\endhead
\bottomrule\noalign{}
\endlastfoot
\(1\) & \\
\(x\) & \\
\(\frac{x^2}{2!} = \frac{x^2}{2}\) & \\
\(\frac{x^3}{3!} = \frac{x^3}{6}\) & \\
\(\frac{x^4}{4!} = \frac{x^4}{24}\) & \\
\(\frac{x^5}{5!} = \frac{x^5}{120}\) & \\
\end{longtable}

\textbf{6.} Write out the derivative series you get:

\(\frac{d}{dx}e^x = 0 + 1 +{}\) \_\_\_\_\_\_ \(+\) \_\_\_\_\_\_ \(+\)
\_\_\_\_\_\_ \(+ \cdots\)

\textbf{7.} Compare to the original series. What do you notice?

\textbf{8.} Complete: \(\frac{d}{dx}e^x ={}\) \_\_\_\_\_\_\_\_\_\_

\textbf{This is remarkable!} The exponential function is its own
derivative!

\begin{center}\rule{0.5\linewidth}{0.5pt}\end{center}

\hypertarget{integrating-ex-term-by-term}{%
\subsubsection{\texorpdfstring{Integrating \(e^x\) Term by
Term}{Integrating e\^{}x Term by Term}}\label{integrating-ex-term-by-term}}

\textbf{9.} Integrate each term of
\(e^x = 1 + x + \frac{x^2}{2} + \frac{x^3}{6} + \frac{x^4}{24} + \cdots\)

\begin{longtable}[]{@{}ll@{}}
\toprule\noalign{}
Term & Integral \\
\midrule\noalign{}
\endhead
\bottomrule\noalign{}
\endlastfoot
\(1\) & \\
\(x\) & \\
\(\frac{x^2}{2}\) & \\
\(\frac{x^3}{6}\) & \\
\(\frac{x^4}{24}\) & \\
\end{longtable}

\textbf{10.} Write out the integrated series:

\(\int e^x \, dx = x + \frac{x^2}{2} +{}\) \_\_\_\_\_\_ \(+\)
\_\_\_\_\_\_ \(+\) \_\_\_\_\_\_ \(+ C\)

\textbf{11.} Factor out from each term to show this equals \(e^x + C\):

Hint: \(x = \frac{x^1}{1!}\), \(\frac{x^2}{2} = \frac{x^2}{2!}\), etc.

\textbf{12.} Complete: \(\int e^x \, dx ={}\) \_\_\_\_\_\_\_\_\_\_
\(+ C\)

\begin{center}\rule{0.5\linewidth}{0.5pt}\end{center}

\hypertarget{the-derivative-and-integral-are-inverses}{%
\subsubsection{The Derivative and Integral are
Inverses}\label{the-derivative-and-integral-are-inverses}}

\textbf{13.} Start with \(f(x) = e^x\). Differentiate, then integrate:

\(f'(x) ={}\) \_\_\_\_\_\_\_\_\_\_

\(\int f'(x) \, dx ={}\) \_\_\_\_\_\_\_\_\_\_

\textbf{14.} Start with \(g(x) = e^x\). Integrate, then differentiate:

\(\int g(x) \, dx ={}\) \_\_\_\_\_\_\_\_\_\_

\(\frac{d}{dx}\left[\int g(x) \, dx\right] ={}\) \_\_\_\_\_\_\_\_\_\_

\begin{center}\rule{0.5\linewidth}{0.5pt}\end{center}

\hypertarget{chain-rule-and-substitution-with-ex}{%
\subsubsection{\texorpdfstring{Chain Rule and Substitution with
\(e^x\)}{Chain Rule and Substitution with e\^{}x}}\label{chain-rule-and-substitution-with-ex}}

\textbf{15.} Using the chain rule:
\(\frac{d}{dx}e^{2x} = e^{2x} \cdot 2 = 2e^{2x}\)

Find: \(\frac{d}{dx}e^{3x} ={}\) \_\_\_\_\_\_\_\_\_\_

\textbf{16.} Find: \(\frac{d}{dx}e^{-x} ={}\) \_\_\_\_\_\_\_\_\_\_

\textbf{17.} Using substitution: \(\int e^{2x} \, dx\)

Let \(u = 2x\), so \(du = 2dx\), meaning \(dx = \frac{1}{2}du\)

\(\int e^{2x} \, dx = \int e^u \cdot \frac{1}{2} du = \frac{1}{2}e^u + C ={}\)
\_\_\_\_\_\_\_\_\_\_ \(+ C\)

\textbf{18.} Find: \(\int e^{5x} \, dx ={}\) \_\_\_\_\_\_\_\_\_\_
\(+ C\)

\begin{center}\rule{0.5\linewidth}{0.5pt}\end{center}

\hypertarget{part-b-sine-and-cosine-via-series-24-problems}{%
\subsection{Part B: Sine and Cosine via Series (24
problems)}\label{part-b-sine-and-cosine-via-series-24-problems}}

\hypertarget{the-series-definitions}{%
\subsubsection{The Series Definitions}\label{the-series-definitions}}

\[\boxed{\sin(x) = x - \frac{x^3}{3!} + \frac{x^5}{5!} - \frac{x^7}{7!} + \cdots = \sum_{n=0}^{\infty} \frac{(-1)^n x^{2n+1}}{(2n+1)!}}\]

\[\boxed{\cos(x) = 1 - \frac{x^2}{2!} + \frac{x^4}{4!} - \frac{x^6}{6!} + \cdots = \sum_{n=0}^{\infty} \frac{(-1)^n x^{2n}}{(2n)!}}\]

\textbf{Note:} \(x\) is in \textbf{radians}, not degrees!

\begin{center}\rule{0.5\linewidth}{0.5pt}\end{center}

\textbf{Understanding the Patterns:}

\textbf{19.} The sine series has only \_\_\_\_\_\_\_\_\_\_ powers of
\(x\) (odd/even).

\textbf{20.} The cosine series has only \_\_\_\_\_\_\_\_\_\_ powers of
\(x\) (odd/even).

\textbf{21.} Both series alternate signs: \(+ \, - \, + \, - \, \cdots\)

\textbf{22.} Use the first 3 terms to approximate \(\sin(0.5)\):

\(\sin(0.5) \approx 0.5 - \frac{(0.5)^3}{6} + \frac{(0.5)^5}{120} ={}\)
\_\_\_\_\_\_\_\_\_\_

Actual: \(\sin(0.5) \approx 0.4794\). How close?

\textbf{23.} Use the first 3 terms to approximate \(\cos(0.5)\):

\(\cos(0.5) \approx 1 - \frac{(0.5)^2}{2} + \frac{(0.5)^4}{24} ={}\)
\_\_\_\_\_\_\_\_\_\_

Actual: \(\cos(0.5) \approx 0.8776\). How close?

\begin{center}\rule{0.5\linewidth}{0.5pt}\end{center}

\hypertarget{discovering-the-derivative-of-sine}{%
\subsubsection{Discovering the Derivative of
Sine}\label{discovering-the-derivative-of-sine}}

\textbf{24.} Differentiate each term of
\(\sin(x) = x - \frac{x^3}{6} + \frac{x^5}{120} - \frac{x^7}{5040} + \cdots\)

\begin{longtable}[]{@{}ll@{}}
\toprule\noalign{}
Term & Derivative \\
\midrule\noalign{}
\endhead
\bottomrule\noalign{}
\endlastfoot
\(x\) & \\
\(-\frac{x^3}{6}\) & \\
\(\frac{x^5}{120}\) & \\
\(-\frac{x^7}{5040}\) & \\
\end{longtable}

\textbf{25.} Write the derivative series:

\(\frac{d}{dx}\sin(x) = 1 -{}\) \_\_\_\_\_\_ \(+\) \_\_\_\_\_\_ \(-\)
\_\_\_\_\_\_ \(+ \cdots\)

\textbf{26.} Compare to the cosine series. What is
\(\frac{d}{dx}\sin(x)\)?

\[\boxed{\frac{d}{dx}\sin(x) = \cos(x)}\]

\begin{center}\rule{0.5\linewidth}{0.5pt}\end{center}

\hypertarget{discovering-the-derivative-of-cosine}{%
\subsubsection{Discovering the Derivative of
Cosine}\label{discovering-the-derivative-of-cosine}}

\textbf{27.} Differentiate each term of
\(\cos(x) = 1 - \frac{x^2}{2} + \frac{x^4}{24} - \frac{x^6}{720} + \cdots\)

\begin{longtable}[]{@{}ll@{}}
\toprule\noalign{}
Term & Derivative \\
\midrule\noalign{}
\endhead
\bottomrule\noalign{}
\endlastfoot
\(1\) & \\
\(-\frac{x^2}{2}\) & \\
\(\frac{x^4}{24}\) & \\
\(-\frac{x^6}{720}\) & \\
\end{longtable}

\textbf{28.} Write the derivative series:

\(\frac{d}{dx}\cos(x) = 0 - x +{}\) \_\_\_\_\_\_ \(-\) \_\_\_\_\_\_
\(+ \cdots\)

\textbf{29.} Factor out \(-1\) and compare to the sine series:

\(\frac{d}{dx}\cos(x) = -\left(x - \frac{x^3}{6} + \frac{x^5}{120} - \cdots\right) ={}\)
\_\_\_\_\_\_\_\_\_\_

\[\boxed{\frac{d}{dx}\cos(x) = -\sin(x)}\]

\begin{center}\rule{0.5\linewidth}{0.5pt}\end{center}

\hypertarget{the-beautiful-cycle}{%
\subsubsection{The Beautiful Cycle}\label{the-beautiful-cycle}}

\textbf{30.} Fill in the differentiation cycle:

\(\sin(x) \xrightarrow{\frac{d}{dx}}\) \_\_\_\_\_\_
\(\xrightarrow{\frac{d}{dx}}\) \_\_\_\_\_\_
\(\xrightarrow{\frac{d}{dx}}\) \_\_\_\_\_\_
\(\xrightarrow{\frac{d}{dx}}\) \_\_\_\_\_\_

\textbf{31.} How many derivatives does it take to get back to
\(\sin(x)\)?

\textbf{32.} What is \(\frac{d^4}{dx^4}\sin(x)\)? (The fourth
derivative)

\textbf{33.} What is \(\frac{d^{100}}{dx^{100}}\cos(x)\)?

\emph{Hint: \(100 = 4 \times 25\), so we go around the cycle 25 times.}

\begin{center}\rule{0.5\linewidth}{0.5pt}\end{center}

\hypertarget{integrating-sine-and-cosine}{%
\subsubsection{Integrating Sine and
Cosine}\label{integrating-sine-and-cosine}}

\textbf{34.} Since \(\frac{d}{dx}\sin(x) = \cos(x)\), what is
\(\int \cos(x) \, dx\)?

\textbf{35.} Since \(\frac{d}{dx}\cos(x) = -\sin(x)\), what is
\(\int \sin(x) \, dx\)?

\emph{Careful with the sign!}

\textbf{36.} Verify by differentiating: \(\frac{d}{dx}(-\cos(x)) ={}\)
\_\_\_\_\_\_\_\_\_\_

\textbf{37.} Fill in the integration cycle:

\(\sin(x) \xrightarrow{\int}\) \_\_\_\_\_\_ \(\xrightarrow{\int}\)
\_\_\_\_\_\_ \(\xrightarrow{\int}\) \_\_\_\_\_\_ \(\xrightarrow{\int}\)
\_\_\_\_\_\_

\emph{(Don't forget the \(+C\) in practice, but ignore it for the
cycle)}

\begin{center}\rule{0.5\linewidth}{0.5pt}\end{center}

\hypertarget{chain-rule-and-substitution-with-trig}{%
\subsubsection{Chain Rule and Substitution with
Trig}\label{chain-rule-and-substitution-with-trig}}

\textbf{38.} Find: \(\frac{d}{dx}\sin(2x) ={}\) \_\_\_\_\_\_\_\_\_\_

\textbf{39.} Find: \(\frac{d}{dx}\cos(3x) ={}\) \_\_\_\_\_\_\_\_\_\_

\textbf{40.} Find: \(\frac{d}{dx}\sin(x^2) ={}\) \_\_\_\_\_\_\_\_\_\_

\textbf{41.} Find: \(\int \cos(4x) \, dx ={}\) \_\_\_\_\_\_\_\_\_\_
\(+ C\)

\textbf{42.} Find: \(\int \sin(2x) \, dx ={}\) \_\_\_\_\_\_\_\_\_\_
\(+ C\)

\begin{center}\rule{0.5\linewidth}{0.5pt}\end{center}

\hypertarget{part-c-the-natural-logarithm-lnx-16-problems}{%
\subsection{\texorpdfstring{Part C: The Natural Logarithm \(\ln(x)\) (16
problems)}{Part C: The Natural Logarithm \textbackslash ln(x) (16 problems)}}\label{part-c-the-natural-logarithm-lnx-16-problems}}

\hypertarget{a-different-approach}{%
\subsubsection{A Different Approach}\label{a-different-approach}}

Unlike \(e^x\), \(\sin\), and \(\cos\), the series for \(\ln(x)\)
doesn't converge for all \(x\). But we can use a related series:

\[\ln(1 + x) = x - \frac{x^2}{2} + \frac{x^3}{3} - \frac{x^4}{4} + \cdots \quad \text{for } -1 < x \leq 1\]

\begin{center}\rule{0.5\linewidth}{0.5pt}\end{center}

\textbf{43.} Use the first 4 terms to approximate
\(\ln(1.5) = \ln(1 + 0.5)\):

\(\ln(1.5) \approx 0.5 - \frac{0.25}{2} + \frac{0.125}{3} - \frac{0.0625}{4} ={}\)
\_\_\_\_\_\_\_\_\_\_

Actual: \(\ln(1.5) \approx 0.4055\). How close?

\textbf{44.} Use the first 4 terms to approximate
\(\ln(2) = \ln(1 + 1)\):

\(\ln(2) \approx 1 - \frac{1}{2} + \frac{1}{3} - \frac{1}{4} ={}\)
\_\_\_\_\_\_\_\_\_\_

Actual: \(\ln(2) \approx 0.6931\). (The series converges slowly here!)

\begin{center}\rule{0.5\linewidth}{0.5pt}\end{center}

\hypertarget{the-derivative-of-lnx}{%
\subsubsection{\texorpdfstring{The Derivative of
\(\ln(x)\)}{The Derivative of \textbackslash ln(x)}}\label{the-derivative-of-lnx}}

\textbf{45.} Differentiate each term of
\(\ln(1+x) = x - \frac{x^2}{2} + \frac{x^3}{3} - \frac{x^4}{4} + \cdots\)

\begin{longtable}[]{@{}ll@{}}
\toprule\noalign{}
Term & Derivative \\
\midrule\noalign{}
\endhead
\bottomrule\noalign{}
\endlastfoot
\(x\) & \\
\(-\frac{x^2}{2}\) & \\
\(\frac{x^3}{3}\) & \\
\(-\frac{x^4}{4}\) & \\
\end{longtable}

\textbf{46.} Write the derivative series:

\(\frac{d}{dx}\ln(1+x) = 1 - x +{}\) \_\_\_\_\_\_ \(-\) \_\_\_\_\_\_
\(+ \cdots\)

\textbf{47.} This is a geometric series! Recall:
\(1 - x + x^2 - x^3 + \cdots = \frac{1}{1+x}\)

So: \(\frac{d}{dx}\ln(1+x) ={}\) \_\_\_\_\_\_\_\_\_\_

\textbf{48.} Using the chain rule: If \(u = 1 + x\), then
\(\frac{d}{dx}\ln(u) = \frac{1}{u} \cdot \frac{du}{dx} = \frac{1}{1+x} \cdot 1\)

For general \(\ln(x)\): \(\frac{d}{dx}\ln(x) ={}\) \_\_\_\_\_\_\_\_\_\_

\[\boxed{\frac{d}{dx}\ln(x) = \frac{1}{x}}\]

\begin{center}\rule{0.5\linewidth}{0.5pt}\end{center}

\hypertarget{the-integral-of-frac1x}{%
\subsubsection{\texorpdfstring{The Integral of
\(\frac{1}{x}\)}{The Integral of \textbackslash frac\{1\}\{x\}}}\label{the-integral-of-frac1x}}

\textbf{49.} We now know: \(\int \frac{1}{x} \, dx = \ln|x| + C\)

Why the absolute value? Because \(\ln(x)\) is only defined for
\(x > 0\), but \(\frac{1}{x}\) exists for \(x \neq 0\).

\textbf{50.} Verify: \(\frac{d}{dx}\ln|x| = \frac{1}{x}\) for both
\(x > 0\) and \(x < 0\).

\begin{center}\rule{0.5\linewidth}{0.5pt}\end{center}

\hypertarget{chain-rule-and-substitution-with-ln}{%
\subsubsection{\texorpdfstring{Chain Rule and Substitution with
\(\ln\)}{Chain Rule and Substitution with \textbackslash ln}}\label{chain-rule-and-substitution-with-ln}}

\textbf{51.} Find: \(\frac{d}{dx}\ln(2x) ={}\) \_\_\_\_\_\_\_\_\_\_

\textbf{52.} Find: \(\frac{d}{dx}\ln(x^2) ={}\) \_\_\_\_\_\_\_\_\_\_

\emph{Two ways: chain rule, or use \(\ln(x^2) = 2\ln(x)\) first}

\textbf{53.} Find: \(\frac{d}{dx}\ln(x^2 + 1) ={}\) \_\_\_\_\_\_\_\_\_\_

\textbf{54.} Find: \(\int \frac{2x}{x^2 + 1} \, dx ={}\)
\_\_\_\_\_\_\_\_\_\_ \(+ C\)

\emph{Hint: Let \(u = x^2 + 1\)}

\textbf{55.} Find: \(\int \frac{3x^2}{x^3 + 5} \, dx ={}\)
\_\_\_\_\_\_\_\_\_\_ \(+ C\)

\textbf{56.} Find: \(\int \frac{1}{2x + 3} \, dx ={}\)
\_\_\_\_\_\_\_\_\_\_ \(+ C\)

\begin{center}\rule{0.5\linewidth}{0.5pt}\end{center}

\hypertarget{the-ln-and-ex-connection}{%
\subsubsection{\texorpdfstring{The \(\ln\) and \(e^x\)
Connection}{The \textbackslash ln and e\^{}x Connection}}\label{the-ln-and-ex-connection}}

\textbf{57.} Since \(e^x\) and \(\ln(x)\) are inverse functions:

\(\ln(e^x) ={}\) \_\_\_\_\_\_\_\_\_\_ and \(e^{\ln(x)} ={}\)
\_\_\_\_\_\_\_\_\_\_

\textbf{58.} Differentiate \(\ln(e^x)\) using the chain rule and verify
it equals 1.

\begin{center}\rule{0.5\linewidth}{0.5pt}\end{center}

\hypertarget{part-d-putting-it-all-together-18-problems}{%
\subsection{Part D: Putting It All Together (18
problems)}\label{part-d-putting-it-all-together-18-problems}}

\hypertarget{summary-table}{%
\subsubsection{Summary Table}\label{summary-table}}

\begin{longtable}[]{@{}lll@{}}
\toprule\noalign{}
Function & Derivative & Integral \\
\midrule\noalign{}
\endhead
\bottomrule\noalign{}
\endlastfoot
\(e^x\) & \(e^x\) & \(e^x + C\) \\
\(\sin(x)\) & \(\cos(x)\) & \(-\cos(x) + C\) \\
\(\cos(x)\) & \(-\sin(x)\) & \(\sin(x) + C\) \\
\(\ln(x)\) & \(\frac{1}{x}\) & \(x\ln(x) - x + C\) \\
\(\frac{1}{x}\) & \(-\frac{1}{x^2}\) & \(\ln|x| + C\) \\
\end{longtable}

\begin{center}\rule{0.5\linewidth}{0.5pt}\end{center}

\hypertarget{mixed-practice-derivatives}{%
\subsubsection{Mixed Practice ---
Derivatives}\label{mixed-practice-derivatives}}

\textbf{59.} \(\frac{d}{dx}(e^x + \sin(x))\)

\textbf{60.} \(\frac{d}{dx}(x^2 + \cos(x))\)

\textbf{61.} \(\frac{d}{dx}(e^x \cdot x)\) --- \emph{Product rule!}

\textbf{62.} \(\frac{d}{dx}(x \cdot \sin(x))\)

\textbf{63.} \(\frac{d}{dx}\left(\frac{e^x}{x}\right)\) ---
\emph{Quotient rule!}

\textbf{64.} \(\frac{d}{dx}(e^{\sin(x)})\) --- \emph{Chain rule!}

\textbf{65.} \(\frac{d}{dx}(\sin(e^x))\)

\textbf{66.} \(\frac{d}{dx}(\ln(\cos(x)))\)

\begin{center}\rule{0.5\linewidth}{0.5pt}\end{center}

\hypertarget{mixed-practice-integrals}{%
\subsubsection{Mixed Practice ---
Integrals}\label{mixed-practice-integrals}}

\textbf{67.} \(\int (e^x + x^2) \, dx\)

\textbf{68.} \(\int (3\cos(x) - 2\sin(x)) \, dx\)

\textbf{69.} \(\int \frac{e^x + 1}{e^x} \, dx\) --- \emph{Simplify
first!}

\textbf{70.} \(\int x \cdot e^x \, dx\) --- \emph{Integration by parts!}

Let \(u = x\), \(dv = e^x dx\). Then \(du = dx\), \(v = e^x\).

\(= xe^x - \int e^x \, dx ={}\) \_\_\_\_\_\_\_\_\_\_

\textbf{71.} \(\int x \cdot \cos(x) \, dx\) --- \emph{Integration by
parts!}

\textbf{72.} \(\int e^x \cdot \sin(x) \, dx\) --- \emph{This requires
parts twice!} (Challenge)

\begin{center}\rule{0.5\linewidth}{0.5pt}\end{center}

\hypertarget{eulers-formula-the-grand-unification}{%
\subsubsection{Euler's Formula --- The Grand
Unification}\label{eulers-formula-the-grand-unification}}

\textbf{73.} Recall: \(e^{ix} = \cos(x) + i\sin(x)\) where
\(i = \sqrt{-1}\)

Using the series for \(e^x\) with \(x\) replaced by \(ix\):

\(e^{ix} = 1 + (ix) + \frac{(ix)^2}{2!} + \frac{(ix)^3}{3!} + \frac{(ix)^4}{4!} + \cdots\)

Calculate the first few powers of \(i\): - \(i^1 = i\) - \(i^2 = -1\) -
\(i^3 = -i\) - \(i^4 = 1\) - \(i^5 ={}\) \_\_\_\_\_

\textbf{74.} Separate the real and imaginary parts of \(e^{ix}\):

Real part: \(1 - \frac{x^2}{2!} + \frac{x^4}{4!} - \cdots ={}\)
\_\_\_\_\_\_\_\_\_\_

Imaginary part: \(x - \frac{x^3}{3!} + \frac{x^5}{5!} - \cdots ={}\)
\_\_\_\_\_\_\_\_\_\_

\textbf{75.} This proves: \(e^{ix} = \cos(x) + i\sin(x)\)

At \(x = \pi\): \(e^{i\pi} = \cos(\pi) + i\sin(\pi) = -1 + 0 = -1\)

So: \(e^{i\pi} + 1 = 0\)

This is called \textbf{Euler's identity} --- often called the most
beautiful equation in mathematics!

\textbf{76.} Explain in your own words why \(e\), \(i\), \(\pi\), \(1\),
and \(0\) appearing in one equation is remarkable.

\begin{center}\rule{0.5\linewidth}{0.5pt}\end{center}

\end{document}
